\chapter{Команды воспроизведения MVP}
\label{appendix:runbook}

Ниже приведен минимальный порядок запуска базового сценария. Предполагается, что зависимости Python установлены, Qdrant доступен по адресу \texttt{http://localhost:6333}, а локальная LLM доступна через OpenAI-compatible endpoint.

\begin{lstlisting}[language=bash,caption={Запуск Qdrant}]
docker compose up -d qdrant
\end{lstlisting}

\begin{lstlisting}[language=bash,caption={Подготовка корпуса и индекса}]
python scripts/generate_documents_registry.py
python scripts/parse_documents.py
python scripts/validate_parsed_corpus.py
python scripts/chunk_documents.py
python scripts/embed_chunks.py
python scripts/index_qdrant.py
\end{lstlisting}

\begin{lstlisting}[language=bash,caption={Проверка поиска и генерации ответа}]
python scripts/search_qdrant.py "какие требования к Wi-Fi сети?" --top-k 3
python scripts/answer_question.py "какие требования к Wi-Fi сети?" --embedding-device cpu
\end{lstlisting}

Для первой MVP-версии центральным является dense search по Qdrant. Дополнительные режимы поиска и повторного ранжирования не включаются в основной сценарий данной записки.

%%% Local Variables:
%%% TeX-engine: xetex
%%% End:
